% Young 表講義第 1 段：第 1 章回顧與第 2 章（動畫 F01–F08、F09–F13）
% 圖與表由 fulton/make_figs.py 產生（含數字檢查）；格子 (i,j) = 第 i 列第 j 行，1 起算。
% figs.tex —— 由 fulton/make_figs.py 產生，請勿手動修改。
% 重新產生：python fulton/make_figs.py（附錄 B，動畫 F09–F19）

\newcommand{\FulFactor}{%
\begin{tabular}{@{}l@{\qquad}l@{}}$V=\begin{ytableau}1 & 2 & 3 \\ 4 & 5 \\ 6\end{ytableau}=$ & $\begin{ytableau}1 & 4 \\ 6\end{ytableau}\cdot\begin{ytableau}2 & 3 \\ 5\end{ytableau}$\quad 或\quad $\begin{ytableau}1 & 6 \\ 4\end{ytableau}\cdot\begin{ytableau}2 & 3 \\ 5\end{ytableau}$\\[10pt]$U_0=\begin{ytableau}1 & 1 \\ 2\end{ytableau}$ & $\begin{ytableau}*(gray!30) {} & *(gray!30) {} & 1 \\ *(gray!30) {} & 1 \\ 2\end{ytableau}\ \to\ \begin{ytableau}1 & 1 \\ 2\end{ytableau}$\ \textcolor{cC}{\checkmark}\qquad $\begin{ytableau}*(gray!30) {} & *(gray!30) {} & 1 \\ *(gray!30) {} & 2 \\ 1\end{ytableau}\ \to\ \begin{ytableau}1 & 1 \\ 2\end{ytableau}$\ \textcolor{cC}{\checkmark}\qquad $\begin{ytableau}*(gray!30) {} & *(gray!30) {} & 2 \\ *(gray!30) {} & 1 \\ 1\end{ytableau}\ \to\ \begin{ytableau}1 & 1 & 2\end{ytableau}$\ \textcolor{cBad}{$\times$}\end{tabular}}

\newcommand{\FulLRAll}{%
{\ytableausetup{smalltableaux}\begin{tabular}{@{}r@{\quad}l@{}}$c=1$ & \begin{tabular}{c}$\begin{ytableau}*(gray!30) {} & *(gray!30) {} & *(gray!30) {} & *(cLeft!35) 1 & *(cLeft!35) 1 \\ *(gray!30) {} & *(gray!30) {} & *(gray!30) {} & *(cB!40) 2 \\ *(gray!30) {} & *(cLeft!35) 1 & *(cLeft!35) 1 \\ *(cLeft!35) 1 & *(cB!40) 2\end{ytableau}$\\[-1pt]{\scriptsize $(5,2)$}\end{tabular}\ \ \begin{tabular}{c}$\begin{ytableau}*(gray!30) {} & *(gray!30) {} & *(gray!30) {} & *(cLeft!35) 1 & *(cLeft!35) 1 \\ *(gray!30) {} & *(gray!30) {} & *(gray!30) {} & *(cB!40) 2 \\ *(gray!30) {} & *(cLeft!35) 1 & *(cLeft!35) 1 \\ *(cLeft!35) 1 & *(cRight!30) 3\end{ytableau}$\\[-1pt]{\scriptsize $(5,1,1)$}\end{tabular}\ \ \begin{tabular}{c}$\begin{ytableau}*(gray!30) {} & *(gray!30) {} & *(gray!30) {} & *(cLeft!35) 1 & *(cLeft!35) 1 \\ *(gray!30) {} & *(gray!30) {} & *(gray!30) {} & *(cB!40) 2 \\ *(gray!30) {} & *(cLeft!35) 1 & *(cRight!30) 3 \\ *(cLeft!35) 1 & *(cC!30) 4\end{ytableau}$\\[-1pt]{\scriptsize $(4,1,1,1)$}\end{tabular}\ \ \begin{tabular}{c}$\begin{ytableau}*(gray!30) {} & *(gray!30) {} & *(gray!30) {} & *(cLeft!35) 1 & *(cLeft!35) 1 \\ *(gray!30) {} & *(gray!30) {} & *(gray!30) {} & *(cB!40) 2 \\ *(gray!30) {} & *(cB!40) 2 & *(cRight!30) 3 \\ *(cRight!30) 3 & *(cC!30) 4\end{ytableau}$\\[-1pt]{\scriptsize $(2,2,2,1)$}\end{tabular}\\[4pt]
$c=2$ & \begin{tabular}{c}$\begin{ytableau}*(gray!30) {} & *(gray!30) {} & *(gray!30) {} & *(cLeft!35) 1 & *(cLeft!35) 1 \\ *(gray!30) {} & *(gray!30) {} & *(gray!30) {} & *(cB!40) 2 \\ *(gray!30) {} & *(cLeft!35) 1 & *(cLeft!35) 1 \\ *(cB!40) 2 & *(cB!40) 2\end{ytableau}$\,$\begin{ytableau}*(gray!30) {} & *(gray!30) {} & *(gray!30) {} & *(cLeft!35) 1 & *(cLeft!35) 1 \\ *(gray!30) {} & *(gray!30) {} & *(gray!30) {} & *(cB!40) 2 \\ *(gray!30) {} & *(cLeft!35) 1 & *(cB!40) 2 \\ *(cLeft!35) 1 & *(cB!40) 2\end{ytableau}$\\[-1pt]{\scriptsize $(4,3)$}\end{tabular}\ \ \begin{tabular}{c}$\begin{ytableau}*(gray!30) {} & *(gray!30) {} & *(gray!30) {} & *(cLeft!35) 1 & *(cLeft!35) 1 \\ *(gray!30) {} & *(gray!30) {} & *(gray!30) {} & *(cB!40) 2 \\ *(gray!30) {} & *(cLeft!35) 1 & *(cRight!30) 3 \\ *(cB!40) 2 & *(cC!30) 4\end{ytableau}$\,$\begin{ytableau}*(gray!30) {} & *(gray!30) {} & *(gray!30) {} & *(cLeft!35) 1 & *(cLeft!35) 1 \\ *(gray!30) {} & *(gray!30) {} & *(gray!30) {} & *(cB!40) 2 \\ *(gray!30) {} & *(cB!40) 2 & *(cRight!30) 3 \\ *(cLeft!35) 1 & *(cC!30) 4\end{ytableau}$\\[-1pt]{\scriptsize $(3,2,1,1)$}\end{tabular}\\[4pt]
 & \begin{tabular}{c}$\begin{ytableau}*(gray!30) {} & *(gray!30) {} & *(gray!30) {} & *(cLeft!35) 1 & *(cLeft!35) 1 \\ *(gray!30) {} & *(gray!30) {} & *(gray!30) {} & *(cB!40) 2 \\ *(gray!30) {} & *(cLeft!35) 1 & *(cB!40) 2 \\ *(cB!40) 2 & *(cRight!30) 3\end{ytableau}$\,$\begin{ytableau}*(gray!30) {} & *(gray!30) {} & *(gray!30) {} & *(cLeft!35) 1 & *(cLeft!35) 1 \\ *(gray!30) {} & *(gray!30) {} & *(gray!30) {} & *(cB!40) 2 \\ *(gray!30) {} & *(cLeft!35) 1 & *(cRight!30) 3 \\ *(cB!40) 2 & *(cB!40) 2\end{ytableau}$\\[-1pt]{\scriptsize $(3,3,1)$}\end{tabular}\ \ \begin{tabular}{c}$\begin{ytableau}*(gray!30) {} & *(gray!30) {} & *(gray!30) {} & *(cLeft!35) 1 & *(cLeft!35) 1 \\ *(gray!30) {} & *(gray!30) {} & *(gray!30) {} & *(cB!40) 2 \\ *(gray!30) {} & *(cLeft!35) 1 & *(cB!40) 2 \\ *(cRight!30) 3 & *(cRight!30) 3\end{ytableau}$\,$\begin{ytableau}*(gray!30) {} & *(gray!30) {} & *(gray!30) {} & *(cLeft!35) 1 & *(cLeft!35) 1 \\ *(gray!30) {} & *(gray!30) {} & *(gray!30) {} & *(cB!40) 2 \\ *(gray!30) {} & *(cB!40) 2 & *(cRight!30) 3 \\ *(cLeft!35) 1 & *(cRight!30) 3\end{ytableau}$\\[-1pt]{\scriptsize $(3,2,2)$}\end{tabular}\\[4pt]
$c=3$ & \begin{tabular}{c}$\begin{ytableau}*(gray!30) {} & *(gray!30) {} & *(gray!30) {} & *(cLeft!35) 1 & *(cLeft!35) 1 \\ *(gray!30) {} & *(gray!30) {} & *(gray!30) {} & *(cB!40) 2 \\ *(gray!30) {} & *(cLeft!35) 1 & *(cLeft!35) 1 \\ *(cB!40) 2 & *(cRight!30) 3\end{ytableau}$\,$\begin{ytableau}*(gray!30) {} & *(gray!30) {} & *(gray!30) {} & *(cLeft!35) 1 & *(cLeft!35) 1 \\ *(gray!30) {} & *(gray!30) {} & *(gray!30) {} & *(cB!40) 2 \\ *(gray!30) {} & *(cLeft!35) 1 & *(cB!40) 2 \\ *(cLeft!35) 1 & *(cRight!30) 3\end{ytableau}$\,$\begin{ytableau}*(gray!30) {} & *(gray!30) {} & *(gray!30) {} & *(cLeft!35) 1 & *(cLeft!35) 1 \\ *(gray!30) {} & *(gray!30) {} & *(gray!30) {} & *(cB!40) 2 \\ *(gray!30) {} & *(cLeft!35) 1 & *(cRight!30) 3 \\ *(cLeft!35) 1 & *(cB!40) 2\end{ytableau}$\\[-1pt]{\scriptsize $(4,2,1)$}\end{tabular}\\[4pt]\end{tabular}}}

\newcommand{\FulRW}{%
反向編號：$\begin{ytableau}*(gray!30) {} & *(gray!30) {} & *(gray!30) {} & 2 & 1 \\ *(gray!30) {} & *(gray!30) {} & *(gray!30) {} & 3 \\ *(gray!30) {} & 5 & 4 \\ 7 & 6\end{ytableau}$\qquad\qquad$w=1\,1\,2\,3\,1\,2\,1\ \longmapsto\ U(w)=\begin{ytableau}1 & 3 & 6 & 7 \\ 2 & 5 \\ 4\end{ytableau}$}

\newcommand{\FulRWTable}{%
{\footnotesize\setlength{\tabcolsep}{4pt}\begin{tabular}{@{}l|ccccccccc|c@{}}$\mu$ & $(5,2)$ & $(5,1,1)$ & $(4,1,1,1)$ & $(2,2,2,1)$ & $(4,3)$ & $(3,2,1,1)$ & $(3,3,1)$ & $(3,2,2)$ & $(4,2,1)$ & 合計\\\hline$\mu$ 上的標準表 & 14 & 15 & 20 & 14 & 14 & 35 & 21 & 21 & 35 & 189\\ 符合 (i)(ii) 的 & 1 & 1 & 1 & 1 & 2 & 2 & 2 & 2 & 3 & 15\end{tabular}}}

\newcommand{\FulCharFour}{%
\begin{tabular}{@{}l|ccccc@{}}$\lambda\backslash\mu$ & $(4)$ & $(3,1)$ & $(2,2)$ & $(2,1,1)$ & $(1,1,1,1)$\\\hline$\chi^{(4)}$ & 1 & 1 & 1 & 1 & 1\\
$\chi^{(3,1)}$ & -1 & 0 & -1 & 1 & 3\\
$\chi^{(2,2)}$ & 0 & -1 & 2 & 0 & 2\\
$\chi^{(2,1,1)}$ & 1 & 0 & -1 & -1 & 3\\
$\chi^{(1,1,1,1)}$ & -1 & 1 & 1 & -1 & 1\\\midrule $z(\mu)$ & 4 & 3 & 8 & 4 & 24\\\end{tabular}}

\newcommand{\FulRecap}{%
$T=\begin{ytableau}1 & 2 & 2 & *(yellow!45) 3 \\ 2 & 3 & *(yellow!45) 5 & 5 \\ 4 & 4 & *(yellow!45) 6 \\ 5 & 6\end{ytableau}\quad\xrightarrow{\ \leftarrow 2\ }\quad T\leftarrow 2=\begin{ytableau}1 & 2 & 2 & *(yellow!45) 2 \\ 2 & 3 & *(yellow!45) 3 & 5 \\ 4 & 4 & *(yellow!45) 5 \\ 5 & 6 & *(cC!35) 6\end{ytableau}$\qquad $U=\begin{ytableau}1 & 3 \\ 2\end{ytableau}$\qquad $T\cdot U=\begin{ytableau}1 & 1 & 2 & 2 & 3 \\ 2 & 2 & 3 & 5 \\ 3 & 4 & 5 \\ 4 & 6 & 6 \\ 5\end{ytableau}$}

\newcommand{\FulRuns}{%
{\ytableausetup{smalltableaux}$\begin{ytableau}\none & \none & \none & \none & \none & \none & \none & \none & \none & \none & \none & 2 \\ \none & \none & \none & \none & \none & \none & 1 & 2 & 2 & 3 & 3 \\ \none & \none & \none & 2 & 3 & 4 \\ 1 & 3 & 4\end{ytableau}$}}

\newcommand{\FulKnuthRules}{%
\begin{tabular}{@{}c@{\qquad\qquad}c@{}}$\textcolor{cLeft}{K'}:\ y\,z\,x\mapsto y\,x\,z\quad(x<y\le z)$ & $\textcolor{cRight}{K''}:\ x\,z\,y\mapsto z\,x\,y\quad(x\le y<z)$\\[6pt]$\begin{ytableau}y & z\end{ytableau}\leftarrow x\ =\ \begin{ytableau}x & z \\ y\end{ytableau}$ & $\begin{ytableau}x & z\end{ytableau}\leftarrow y\ =\ \begin{ytableau}x & y \\ z\end{ytableau}$\\[4pt]字：$y\,z\,x\ \to\ y\,x\,z$ & 字：$x\,z\,y\ \to\ z\,x\,y$\end{tabular}}

\newcommand{\FulChain}{%
\begin{tabular}{@{}rll@{}}\toprule 步 & 字 & 規則\\\midrule0 & $5\,6\,4\,4\,6\,2\,3\,5\,5\,1\,2\,2\,3\,\textcolor{cB}{\mathbf{2}}$ & $w(T)\cdot 2$\\
1 & $5\,6\,4\,4\,6\,2\,3\,5\,5\,1\,2\,\textcolor{cB}{\mathbf{3}}\,\textcolor{cB}{\mathbf{2}}\,\textcolor{cB}{\mathbf{2}}$ & $\textcolor{cRight}{K''}$\\
2 & $5\,6\,4\,4\,6\,2\,3\,5\,5\,1\,\textcolor{cB}{\mathbf{3}}\,\textcolor{cB}{\mathbf{2}}\,\textcolor{cB}{\mathbf{2}}\,2$ & $\textcolor{cRight}{K''}$\\
3 & $5\,6\,4\,4\,6\,2\,3\,5\,5\,\textcolor{cB}{\mathbf{3}}\,\textcolor{cB}{\mathbf{1}}\,\textcolor{cB}{\mathbf{2}}\,2\,2$ & $\textcolor{cRight}{K''}$\\
4 & $5\,6\,4\,4\,6\,2\,3\,\textcolor{cB}{\mathbf{5}}\,\textcolor{cB}{\mathbf{3}}\,\textcolor{cB}{\mathbf{5}}\,1\,2\,2\,2$ & $\textcolor{cLeft}{K'}$\\
5 & $5\,6\,4\,4\,6\,2\,\textcolor{cB}{\mathbf{5}}\,\textcolor{cB}{\mathbf{3}}\,\textcolor{cB}{\mathbf{3}}\,5\,1\,2\,2\,2$ & $\textcolor{cRight}{K''}$\\
6 & $5\,6\,4\,4\,6\,\textcolor{cB}{\mathbf{5}}\,\textcolor{cB}{\mathbf{2}}\,\textcolor{cB}{\mathbf{3}}\,3\,5\,1\,2\,2\,2$ & $\textcolor{cRight}{K''}$\\
7 & $5\,6\,4\,\textcolor{cB}{\mathbf{6}}\,\textcolor{cB}{\mathbf{4}}\,\textcolor{cB}{\mathbf{5}}\,2\,3\,3\,5\,1\,2\,2\,2$ & $\textcolor{cRight}{K''}$\\
8 & $5\,6\,\textcolor{cB}{\mathbf{6}}\,\textcolor{cB}{\mathbf{4}}\,\textcolor{cB}{\mathbf{4}}\,5\,2\,3\,3\,5\,1\,2\,2\,2$ & $\textcolor{cRight}{K''}$\\\bottomrule\end{tabular}}

\newcommand{\FulSlide}{%
$\begin{ytableau}1 & *(yellow!40) {} & 4 \\ 2 & 3 & 5\end{ytableau}\ \longrightarrow\ \begin{ytableau}1 & 3 & 4 \\ 2 & *(yellow!40) {} & 5\end{ytableau}$\qquad $2\,3\,5\,1\,4\equiv 2\,\textcolor{cB}{\mathbf{3}}\,\textcolor{cB}{\mathbf{1}}\,\textcolor{cB}{\mathbf{5}}\,4\equiv \textcolor{cB}{\mathbf{2}}\,\textcolor{cB}{\mathbf{1}}\,\textcolor{cB}{\mathbf{3}}\,5\,4\equiv 2\,1\,\textcolor{cB}{\mathbf{5}}\,\textcolor{cB}{\mathbf{3}}\,\textcolor{cB}{\mathbf{4}}\equiv \textcolor{cB}{\mathbf{2}}\,\textcolor{cB}{\mathbf{5}}\,\textcolor{cB}{\mathbf{1}}\,3\,4$}

\newcommand{\FulPieri}{%
\begin{tabular}{@{}r@{\ }l@{}}$S_{(2,1)}\cdot S_{(2)}=$ & $\begin{ytableau}{} & {} & *(cLeft!40) {} & *(cLeft!40) {} \\ {}\end{ytableau}+\begin{ytableau}{} & {} & *(cLeft!40) {} \\ {} & *(cLeft!40) {}\end{ytableau}+\begin{ytableau}{} & {} & *(cLeft!40) {} \\ {} \\ *(cLeft!40) {}\end{ytableau}+\begin{ytableau}{} & {} \\ {} & *(cLeft!40) {} \\ *(cLeft!40) {}\end{ytableau}$\\[10pt]$S_{(2,1)}\cdot S_{(1,1)}=$ & $\begin{ytableau}{} & {} & *(cRight!35) {} \\ {} & *(cRight!35) {}\end{ytableau}+\begin{ytableau}{} & {} & *(cRight!35) {} \\ {} \\ *(cRight!35) {}\end{ytableau}+\begin{ytableau}{} & {} \\ {} & *(cRight!35) {} \\ *(cRight!35) {}\end{ytableau}+\begin{ytableau}{} & {} \\ {} \\ *(cRight!35) {} \\ *(cRight!35) {}\end{ytableau}$\end{tabular}}

\newcommand{\FulKostkaTabs}{%
$\begin{ytableau}*(cLeft!35) 1 & *(cLeft!35) 1 & *(cB!40) 2 \\ *(cB!40) 2 & *(cRight!30) 3\end{ytableau}$\ \ {\small $\varnothing\subset (2)\subset (3,1)\subset (3,2)$}\qquad\qquad $\begin{ytableau}*(cLeft!35) 1 & *(cLeft!35) 1 & *(cRight!30) 3 \\ *(cB!40) 2 & *(cB!40) 2\end{ytableau}$\ \ {\small $\varnothing\subset (2)\subset (2,2)\subset (3,2)$}}

\newcommand{\FulKostkaMatrix}{%
\begin{tabular}{@{}l|ccccc@{}}$\lambda\backslash\mu$ & $(1,1,1,1)$ & $(2,1,1)$ & $(2,2)$ & $(3,1)$ & $(4)$\\\hline$(1,1,1,1)$ & \textbf{1} & \textcolor{gray}{0} & \textcolor{gray}{0} & \textcolor{gray}{0} & \textcolor{gray}{0}\\
$(2,1,1)$ & 3 & \textbf{1} & \textcolor{gray}{0} & \textcolor{gray}{0} & \textcolor{gray}{0}\\
$(2,2)$ & 2 & 1 & \textbf{1} & \textcolor{gray}{0} & \textcolor{gray}{0}\\
$(3,1)$ & 3 & 2 & 1 & \textbf{1} & \textcolor{gray}{0}\\
$(4)$ & 1 & 1 & 1 & 1 & \textbf{1}\\\end{tabular}}

\newcommand{\FulCut}{%
$T=\begin{ytableau}*(cLeft!30) 1 & *(cLeft!30) 2 & *(cRight!25) 2 & *(cRight!25) 3 \\ *(cLeft!30) 2 & *(cLeft!30) 3 & *(cRight!25) 5 & *(cRight!25) 5 \\ *(cLeft!30) 4 & *(cLeft!30) 4 & *(cRight!25) 6 \\ *(cLeft!30) 5 & *(cLeft!30) 6\end{ytableau}$\qquad $T'=\begin{ytableau}1 & 2 \\ 2 & 3 \\ 4 & 4 \\ 5 & 6\end{ytableau}$\quad $T''=\begin{ytableau}2 & 3 \\ 5 & 5 \\ 6\end{ytableau}$\qquad$T'*T''=\begin{ytableau}*(gray!25) {} & *(gray!25) {} & 2 & 3 \\ *(gray!25) {} & *(gray!25) {} & 5 & 5 \\ *(gray!25) {} & *(gray!25) {} & 6 \\ 1 & 2 \\ 2 & 3 \\ 4 & 4 \\ 5 & 6\end{ytableau}$}

\newcommand{\FulIncreasing}{%
\begin{tabular}{@{}lll@{}}$k=1$: & $\textcolor{cLeft}{\mathbf{1}}\ 3\ 4\ 2\ 3\ 4\ \textcolor{cLeft}{\mathbf{1}}\ \textcolor{cLeft}{\mathbf{2}}\ \textcolor{cLeft}{\mathbf{2}}\ \textcolor{cLeft}{\mathbf{3}}\ \textcolor{cLeft}{\mathbf{3}}\ 2$ & $L(w,1)=6$\\
$k=2$: & $\textcolor{cB}{\mathbf{1}}\ \textcolor{cB}{\mathbf{3}}\ 4\ 2\ \textcolor{cB}{\mathbf{3}}\ \textcolor{cB}{\mathbf{4}}\ \textcolor{cLeft}{\mathbf{1}}\ \textcolor{cLeft}{\mathbf{2}}\ \textcolor{cLeft}{\mathbf{2}}\ \textcolor{cLeft}{\mathbf{3}}\ \textcolor{cLeft}{\mathbf{3}}\ 2$ & $L(w,2)=9=5+4$\\
$k=3$: & $\textcolor{cLeft}{\mathbf{1}}\ \textcolor{cLeft}{\mathbf{3}}\ \textcolor{cLeft}{\mathbf{4}}\ \textcolor{cB}{\mathbf{2}}\ \textcolor{cB}{\mathbf{3}}\ \textcolor{cLeft}{\mathbf{4}}\ \textcolor{cRight}{\mathbf{1}}\ \textcolor{cRight}{\mathbf{2}}\ \textcolor{cRight}{\mathbf{2}}\ \textcolor{cB}{\mathbf{3}}\ \textcolor{cB}{\mathbf{3}}\ \textcolor{cRight}{\mathbf{2}}$ & $L(w,3)=12=4+4+4$\\\end{tabular}\qquad$P(w)=\begin{ytableau}*(cLeft!30) 1 & *(cLeft!30) 1 & *(cLeft!30) 2 & *(cLeft!30) 2 & *(cLeft!30) 2 & *(cLeft!30) 3 \\ *(cB!30) 2 & *(cB!30) 3 & *(cB!30) 3 \\ *(cRight!30) 3 & *(cRight!30) 4 & *(cRight!30) 4\end{ytableau}$}

\newcommand{\FulRebuild}{%
{\small\begin{tabular}{@{}rlccl@{}}\toprule $p$ & $w_p$（灰色是已去掉的） & $L(w_p,1..3)$ & 形狀 & 最大字母放哪\\\midrule0 & $1\,3\,4\,2\,3\,4\,1\,2\,2\,3\,3\,2$ & 6, 9, 12 & $(6,3,3)$ & $4\to$ 第 3 列\\
1 & $1\,3\,4\,2\,3\,\textcolor{gray!45}{\mathbf{4}}\,1\,2\,2\,3\,3\,2$ & 6, 9, 11 & $(6,3,2)$ & $4\to$ 第 3 列\\
2 & $1\,3\,\textcolor{gray!45}{\mathbf{4}}\,2\,3\,\textcolor{gray!45}{\mathbf{4}}\,1\,2\,2\,3\,3\,2$ & 6, 9, 10 & $(6,3,1)$ & $3\to$ 第 1 列\\
3 & $1\,3\,\textcolor{gray!45}{\mathbf{4}}\,2\,3\,\textcolor{gray!45}{\mathbf{4}}\,1\,2\,2\,3\,\textcolor{gray!45}{\mathbf{3}}\,2$ & 5, 8, 9 & $(5,3,1)$ & $3\to$ 第 2 列\\
4 & $1\,3\,\textcolor{gray!45}{\mathbf{4}}\,2\,3\,\textcolor{gray!45}{\mathbf{4}}\,1\,2\,2\,\textcolor{gray!45}{\mathbf{3}}\,\textcolor{gray!45}{\mathbf{3}}\,2$ & 5, 7, 8 & $(5,2,1)$ & $3\to$ 第 2 列\\
5 & $1\,3\,\textcolor{gray!45}{\mathbf{4}}\,2\,\textcolor{gray!45}{\mathbf{3}}\,\textcolor{gray!45}{\mathbf{4}}\,1\,2\,2\,\textcolor{gray!45}{\mathbf{3}}\,\textcolor{gray!45}{\mathbf{3}}\,2$ & 5, 6, 7 & $(5,1,1)$ & $3\to$ 第 3 列\\
6 & $1\,\textcolor{gray!45}{\mathbf{3}}\,\textcolor{gray!45}{\mathbf{4}}\,2\,\textcolor{gray!45}{\mathbf{3}}\,\textcolor{gray!45}{\mathbf{4}}\,1\,2\,2\,\textcolor{gray!45}{\mathbf{3}}\,\textcolor{gray!45}{\mathbf{3}}\,2$ & 5, 6, 6 & $(5,1)$ & $2\to$ 第 1 列\\
7 & $1\,\textcolor{gray!45}{\mathbf{3}}\,\textcolor{gray!45}{\mathbf{4}}\,2\,\textcolor{gray!45}{\mathbf{3}}\,\textcolor{gray!45}{\mathbf{4}}\,1\,2\,2\,\textcolor{gray!45}{\mathbf{3}}\,\textcolor{gray!45}{\mathbf{3}}\,\textcolor{gray!45}{\mathbf{2}}$ & 4, 5, 5 & $(4,1)$ & $2\to$ 第 1 列\\
8 & $1\,\textcolor{gray!45}{\mathbf{3}}\,\textcolor{gray!45}{\mathbf{4}}\,2\,\textcolor{gray!45}{\mathbf{3}}\,\textcolor{gray!45}{\mathbf{4}}\,1\,2\,\textcolor{gray!45}{\mathbf{2}}\,\textcolor{gray!45}{\mathbf{3}}\,\textcolor{gray!45}{\mathbf{3}}\,\textcolor{gray!45}{\mathbf{2}}$ & 3, 4, 4 & $(3,1)$ & $2\to$ 第 1 列\\
9 & $1\,\textcolor{gray!45}{\mathbf{3}}\,\textcolor{gray!45}{\mathbf{4}}\,2\,\textcolor{gray!45}{\mathbf{3}}\,\textcolor{gray!45}{\mathbf{4}}\,1\,\textcolor{gray!45}{\mathbf{2}}\,\textcolor{gray!45}{\mathbf{2}}\,\textcolor{gray!45}{\mathbf{3}}\,\textcolor{gray!45}{\mathbf{3}}\,\textcolor{gray!45}{\mathbf{2}}$ & 2, 3, 3 & $(2,1)$ & $2\to$ 第 2 列\\
10 & $1\,\textcolor{gray!45}{\mathbf{3}}\,\textcolor{gray!45}{\mathbf{4}}\,\textcolor{gray!45}{\mathbf{2}}\,\textcolor{gray!45}{\mathbf{3}}\,\textcolor{gray!45}{\mathbf{4}}\,1\,\textcolor{gray!45}{\mathbf{2}}\,\textcolor{gray!45}{\mathbf{2}}\,\textcolor{gray!45}{\mathbf{3}}\,\textcolor{gray!45}{\mathbf{3}}\,\textcolor{gray!45}{\mathbf{2}}$ & 2, 2, 2 & $(2)$ & $1\to$ 第 1 列\\
11 & $1\,\textcolor{gray!45}{\mathbf{3}}\,\textcolor{gray!45}{\mathbf{4}}\,\textcolor{gray!45}{\mathbf{2}}\,\textcolor{gray!45}{\mathbf{3}}\,\textcolor{gray!45}{\mathbf{4}}\,\textcolor{gray!45}{\mathbf{1}}\,\textcolor{gray!45}{\mathbf{2}}\,\textcolor{gray!45}{\mathbf{2}}\,\textcolor{gray!45}{\mathbf{3}}\,\textcolor{gray!45}{\mathbf{3}}\,\textcolor{gray!45}{\mathbf{2}}$ & 1, 1, 1 & $(1)$ & $1\to$ 第 1 列\\\bottomrule\end{tabular}}}

\newcommand{\FulRS}{%
{\ytableausetup{smalltableaux}$k=1$：$x_1=5$\ \ $\begin{ytableau}5\end{ytableau}\ \ \begin{ytableau}1\end{ytableau}$\quad $k=2$：$x_2=4$\ \ $\begin{ytableau}4 \\ 5\end{ytableau}\ \ \begin{ytableau}1 \\ 2\end{ytableau}$\quad $k=3$：$x_3=8$\ \ $\begin{ytableau}4 & 8 \\ 5\end{ytableau}\ \ \begin{ytableau}1 & 3 \\ 2\end{ytableau}$\quad $k=4$：$x_4=2$\ \ $\begin{ytableau}2 & 8 \\ 4 \\ 5\end{ytableau}\ \ \begin{ytableau}1 & 3 \\ 2 \\ 4\end{ytableau}$}\\[8pt]$P=\begin{ytableau}1 & 1 & 3 & 5 \\ 2 & 3 \\ 4 & 4 \\ 5 & 7 \\ 8\end{ytableau}\qquad Q=\begin{ytableau}1 & 3 & 6 & 8 \\ 2 & 5 \\ 4 & 9 \\ 7 & 10 \\ 11\end{ytableau}$}

\newcommand{\FulArray}{%
$\omega=\begin{pmatrix}1 & 1 & 1 & 2 & 2 & 3 & 3 & 3 & 3\\1 & 2 & 2 & 1 & 2 & 1 & 1 & 1 & 2\end{pmatrix}\qquad A=\begin{pmatrix}1 & 2\\1 & 1\\3 & 1\end{pmatrix}\qquad P=\begin{ytableau}1 & 1 & 1 & 1 & 1 & 2 \\ 2 & 2 & 2\end{ytableau}\quad Q=\begin{ytableau}1 & 1 & 1 & 2 & 3 & 3 \\ 2 & 3 & 3\end{ytableau}$}

\newcommand{\FulBalls}{%
$A$:\ \begin{tikzpicture}[baseline=(current bounding box.center)]
\draw[gray!60] (0,0.00) -- (2.00,0.00);
\draw[gray!60] (0,-1.00) -- (2.00,-1.00);
\draw[gray!60] (0,-2.00) -- (2.00,-2.00);
\draw[gray!60] (0,-3.00) -- (2.00,-3.00);
\draw[gray!60] (0.00,0) -- (0.00,-3.00);
\draw[gray!60] (1.00,0) -- (1.00,-3.00);
\draw[gray!60] (2.00,0) -- (2.00,-3.00);
\node[font=\scriptsize,gray] at (-0.25,-0.50) {1};
\node[font=\scriptsize,gray] at (-0.25,-1.50) {2};
\node[font=\scriptsize,gray] at (-0.25,-2.50) {3};
\node[font=\scriptsize,gray] at (0.50,0.22) {1};
\node[font=\scriptsize,gray] at (1.50,0.22) {2};
\node[circle,draw=cA,fill=cA!15,inner sep=0pt,minimum size=9pt,font=\tiny] at (0.50,-0.50) {1};
\node[circle,draw=cA,fill=cA!15,inner sep=0pt,minimum size=9pt,font=\tiny] at (1.38,-0.38) {2};
\node[circle,draw=cA,fill=cA!15,inner sep=0pt,minimum size=9pt,font=\tiny] at (1.62,-0.62) {3};
\node[circle,draw=cA,fill=cA!15,inner sep=0pt,minimum size=9pt,font=\tiny] at (0.50,-1.50) {2};
\node[circle,draw=cA,fill=cA!15,inner sep=0pt,minimum size=9pt,font=\tiny] at (1.50,-1.50) {4};
\node[circle,draw=cA,fill=cA!15,inner sep=0pt,minimum size=9pt,font=\tiny] at (0.25,-2.25) {3};
\node[circle,draw=cA,fill=cA!15,inner sep=0pt,minimum size=9pt,font=\tiny] at (0.50,-2.50) {4};
\node[circle,draw=cA,fill=cA!15,inner sep=0pt,minimum size=9pt,font=\tiny] at (0.75,-2.75) {5};
\node[circle,draw=cA,fill=cA!15,inner sep=0pt,minimum size=9pt,font=\tiny] at (1.50,-2.50) {6};\end{tikzpicture}\ \ {\small\begin{tabular}{@{}l@{}}$P$ 第 1 列：$1,1,1,1,1,2$\\$Q$ 第 1 列：$1,1,1,2,3,3$\end{tabular}}\qquad $A^\flat$:\ \begin{tikzpicture}[baseline=(current bounding box.center)]
\draw[gray!60] (0,0.00) -- (2.00,0.00);
\draw[gray!60] (0,-1.00) -- (2.00,-1.00);
\draw[gray!60] (0,-2.00) -- (2.00,-2.00);
\draw[gray!60] (0,-3.00) -- (2.00,-3.00);
\draw[gray!60] (0.00,0) -- (0.00,-3.00);
\draw[gray!60] (1.00,0) -- (1.00,-3.00);
\draw[gray!60] (2.00,0) -- (2.00,-3.00);
\node[font=\scriptsize,gray] at (-0.25,-0.50) {1};
\node[font=\scriptsize,gray] at (-0.25,-1.50) {2};
\node[font=\scriptsize,gray] at (-0.25,-2.50) {3};
\node[font=\scriptsize,gray] at (0.50,0.22) {1};
\node[font=\scriptsize,gray] at (1.50,0.22) {2};
\node[circle,draw=cA,fill=cA!15,inner sep=0pt,minimum size=9pt,font=\tiny] at (1.50,-1.50) {1};
\node[circle,draw=cA,fill=cA!15,inner sep=0pt,minimum size=9pt,font=\tiny] at (1.38,-2.38) {2};
\node[circle,draw=cA,fill=cA!15,inner sep=0pt,minimum size=9pt,font=\tiny] at (1.62,-2.62) {3};\end{tikzpicture}\ \ {\small\begin{tabular}{@{}l@{}}$P$ 第 2 列：$2,2,2$\\$Q$ 第 2 列：$2,3,3$\end{tabular}}}

\newcommand{\FulHook}{%
$h(i,j)=\begin{ytableau}9 & 8 & 7 & 5 & 4 & 1 \\ 7 & 6 & 5 & 3 & 2 \\ 6 & 5 & 4 & 2 & 1 \\ 3 & 2 & 1\end{ytableau}\qquad 5+j-i=\begin{ytableau}*(cB!35) 5 & 6 & 7 & 8 & 9 & 10 \\ 4 & *(cB!35) 5 & 6 & 7 & 8 \\ 3 & 4 & *(cB!35) 5 & 6 & 7 \\ 2 & 3 & 4\end{ytableau}$}


% ---------------------------------------------------------------------
\section{第 1 章回顧：列插入、表的字與滑動（動畫 F01--F08）}\label{sec:f-recap}

這一節只把後面各章要用的定義集中寫一次（動畫 F01--F08 有完整的示範）。

\textbf{Young 圖與表。}分割 $\lambda=(\lambda_1\ge\lambda_2\ge\cdots)$ 畫成 Young 圖：第 $i$ 列有 $\lambda_i$ 格、靠左對齊、第一列在最上面。
在格子裡填正整數，每一列由左到右\textbf{弱遞增}、每一行由上到下\textbf{嚴格遞增}，就是一個（半標準）\textbf{表}；
$1,\ldots,n$ 各用一次的叫\textbf{標準表}。從 $\lambda$ 挖掉左上角的小圖 $\mu$ 得到斜形狀 $\lambda/\mu$，填法相同的叫\textbf{斜表}。

\textbf{列插入 $T\leftarrow x$} \fulref{p.7}。把 $x$ 放進第一列：如果這一列沒有比 $x$ 大的數，$x$ 接在列尾；
否則 $x$ 擠出「最左邊、比 $x$ 大」的那個數，被擠出的數再用同樣的規則插入下一列，直到某個數接在列尾（可能自成新的一列）。
結果仍是表，形狀多出一個外角。知道新格子在哪裡，就能反向擠出，還原 $T$ 與 $x$。

\textbf{表的字與乘積} \fulref{p.11}。表的\textbf{字} $w(T)$：從最下面一列讀到最上面一列，每列由左往右。
乘積 $T\cdot U$：把 $w(U)$ 的字母依序列插入 $T$。

\begin{figure}[H]
\centering
\FulRecap
\caption{列插入 $T\leftarrow 2$：2 擠出 3、3 擠出 5、5 擠出 6、6 接在第 4 列尾（\colorbox{yellow!45}{黃色}是擠壓路徑，
\colorbox{cC!35}{綠色}是新格子 $(4,3)$）。右邊是乘積 $T\cdot U$：把 $w(U)=2\,1\,3$ 依序插入 $T$。
這裡 $w(T)=5\,6\,4\,4\,6\,2\,3\,5\,5\,1\,2\,2\,3$。}
\label{fig:f-recap}
\end{figure}

\textbf{滑動與 rectification} \fulref{p.12--13}。在斜表裡挑一個\textbf{內角}（$\mu$ 的角）當洞，
把洞右邊與下面兩個數中\textbf{較小}的滑進洞（相等時取下面的），洞跟著移動，直到洞成為外角、被拿掉。
一直滑到沒有內角，得到一個表，叫做 $S$ 的 \textbf{rectification} $\mathrm{Rect}(S)$。
把 $T$ 放在左下、$U$ 放在右上、中間補一塊空白矩形，得到斜表 $T*U$。

\begin{keybox}[第 1 章留下的三個 Claim（第 2--3 章證明它們）]
\begin{enumerate}[nosep]
\item 乘積有結合律：$(T\cdot U)\cdot V=T\cdot(U\cdot V)$。
\item $\mathrm{Rect}(S)$ 與內角的挑法無關。
\item $\mathrm{Rect}(T*U)=T\cdot U$。
\end{enumerate}
三者都由第 2 章的主定理（\S\ref{sec:f-slides}）推出，主定理的唯一性在第 3 章（\S\ref{sec:f-unique}）證明。
\end{keybox}

% ---------------------------------------------------------------------
\section{第 2 章：字與 plactic monoid（動畫 F09--F13）}\label{sec:f-ch2}

\subsection{字與 Knuth 變換}\label{sec:f-knuth}

\textbf{由字還原表} \fulref{p.17}。把字在每個「下一個字母比前一個小」的地方切開，切出的弱遞增片段由下往上疊起來，
就是原來的表。例如 $5\,6\,4\,4\,6\,2\,3\,5\,5\,1\,2\,2\,3$ 切成 $56\mid 446\mid 2355\mid 1223$，正是圖~\ref{fig:f-recap} 中 $T$ 的四列。
不是每個字都來自表：$3\,4\,1\,4$ 切成 $34\mid 14$，疊起來第 2 行是 4 在 4 上面，不是嚴格遞增。
但\textbf{每個字都是某個斜表的字}：每段當一列，每一列放在下一列的整個右上方，各行至多一格，條件自動成立（圖~\ref{fig:f-runs}）。

\begin{figure}[H]
\centering
\FulRuns
\caption{字 $w=1\,3\,4\,2\,3\,4\,1\,2\,2\,3\,3\,2$ 切成 $134\mid 234\mid 12233\mid 2$，排成階梯形的斜表，它的字就是 $w$。}
\label{fig:f-runs}
\end{figure}

\textbf{用字的語言寫列插入} \fulref{p.18}。把一列的字寫成 $u\,x'\,v$：$x'$ 是第一個比 $x$ 大的字母，
所以 $u\le x<x'\le v$（意思是 $u$ 的每個字母 $\le x$、$v$ 的每個字母 $\ge x'$）。插入 $x$ 就是
\[
(u\,x'\,v)\cdot x\ \longmapsto\ x'\cdot(u\,x\,v),
\]
被擠出的 $x'$ 移到這一列的左邊，而在字裡，下一列恰好就在左邊。圖~\ref{fig:f-recap} 的插入寫成四行：
\[
\begin{aligned}
(56)(446)(2355)(1223)\cdot 2&\mapsto(56)(446)(2355)\cdot 3\cdot(1222)\mapsto(56)(446)\cdot 5\cdot(2335)(1222)\\
&\mapsto(56)\cdot 6\cdot(445)(2335)(1222)=(566)(445)(2335)(1222).
\end{aligned}
\]

\textbf{Knuth 的拆解} \fulref{p.18--19}。每一行還可以再拆成更小的步：先讓 $x$ 一格一格往左移過 $v$ 的字母，
再讓 $x'$ 一格一格往左移過 $u$ 的字母。每一步只動三個相鄰的字母、交換其中兩個，而且只有兩種：
\begin{figure}[H]
\centering
\FulKnuthRules
\caption{兩種基本變換與記法：$K'$ 是把 $x$ 插入一列 $[\,y\ z\,]$（$x$ 擠出 $y$），$K''$ 是把 $y$ 插入 $[\,x\ z\,]$（$y$ 擠出 $z$）。
插入前後的字分別是左邊與右邊的三個字母。}
\label{fig:f-knuth-rules}
\end{figure}
換句話說：一個字母 $y$ 同一側的兩個鄰居若一個比 $y$ 小、一個比 $y$ 大，就可以交換它們
（有相等時，照「相等的字母左邊的算小」來判斷）。

\begin{keybox}[Knuth 等價]
對字裡三個相鄰的字母做一次 $K'$、$K''$ 或它們的反方向，叫做一次\textbf{基本 Knuth 變換}。
兩個字若能用一串基本變換互相變到，就說它們 \textbf{Knuth 等價}，記作 $w\equiv w'$。
\end{keybox}

\begin{thmbox}[命題 1 \fulref{p.19}]
對任何表 $T$ 與正整數 $x$，$w(T\leftarrow x)\equiv w(T)\cdot x$。所以 $w(T\cdot U)\equiv w(T)\cdot w(U)$。
\end{thmbox}

表~\ref{tab:f-chain} 是圖~\ref{fig:f-recap} 的插入拆成的 8 步。第 1 列 $1\,2\,2\,3$：$x=2$ 擠出 $x'=3$，$v$ 是空的，
$3$ 往左移過 $u=1\,2\,2$，三步 $K''$；第 2 列：$3$ 擠出 $5$，先用一步 $K'$ 越過 $v=5$，再用兩步 $K''$ 越過 $u=2\,3$；
第 3 列：$5$ 擠出 $6$，兩步 $K''$ 越過 $4\,4$；第 4 列：$6$ 接在 $5\,6$ 後面，字已經是 $w(T\leftarrow 2)$。

\begin{table}[H]
\centering
\FulChain
\caption{$w(T)\cdot 2\equiv w(T\leftarrow 2)$ 的 8 步基本變換（每一步動到的三個字母標成橘色）。}
\label{tab:f-chain}
\end{table}

\begin{notebox}
一個等價類有多少個字？若 $T$ 的形狀是 $\lambda$，與 $w(T)$ 等價的字恰有 $f^\lambda$ 個（$f^\lambda$ = 形狀 $\lambda$ 的標準表個數，
書 p.57 練習 16，理由見 \S\ref{sec:f-rs}）。例如 $1\,3\,2\equiv 3\,1\,2$，而 $f^{(2,1)}=2$。
\end{notebox}

% ---------------------------------------------------------------------
\subsection{滑動、主定理與 plactic monoid}\label{sec:f-slides}

滑動的中途會出現「有洞」的表；它的字照常讀，跳過洞。\textbf{水平滑動}（洞右邊的數往左移）字完全不變；
\textbf{垂直滑動}（洞下面的數往上移）字會變，但仍然 Knuth 等價。最簡單的情形 \fulref{p.20}：上列 $u\ \square\ y$、
下列 $v\ x\ z$，$x$ 往上滑。用 $u,v,x,y,z=1,2,3,4,5$，字從 $v\,x\,z\,u\,y=2\,3\,5\,1\,4$ 變成 $v\,z\,u\,x\,y=2\,5\,1\,3\,4$，
四步基本變換就能接起來（圖~\ref{fig:f-slide}）。一般情形把 $u,v,y,z$ 換成整段的遞增列，對列長用歸納法，每一步用一次命題 1 \fulref{p.21--22}。

\begin{figure}[H]
\centering
\FulSlide
\caption{垂直滑動前後的字 Knuth 等價：$23514\equiv 25134$（每一步動到的三個字母標成橘色）。}
\label{fig:f-slide}
\end{figure}

\begin{thmbox}[命題 2 \fulref{p.22}]
若一個斜表可以經由一串滑動變成另一個，它們的字 Knuth 等價。
\end{thmbox}

\begin{thmbox}[主定理 \fulref{p.22}]
每個字都恰好與\textbf{一個}表的字 Knuth 等價。
\end{thmbox}

\textbf{存在}很容易：把 $w=x_1x_2\cdots x_r$ 的字母依序插入空表，
\[
P(w)=\bigl(\cdots\bigl((x_1\leftarrow x_2)\leftarrow x_3\bigr)\leftarrow\cdots\bigr)\leftarrow x_r ,
\]
由命題 1，每插入一個字母，等價類都不變，所以 $w\equiv w(P(w))$。例如 $w=1\,3\,4\,2\,3\,4\,1\,2\,2\,3\,3\,2$ 的
$P(w)$ 是圖~\ref{fig:f-increasing} 右邊的表，它的字 $3\,4\,4\,2\,3\,3\,1\,1\,2\,2\,2\,3$ 和 $w$ 不同，卻彼此等價。
\textbf{唯一}則需要新的想法，留到第 3 章（\S\ref{sec:f-unique}）。先承認它，就得到：

\begin{thmbox}[推論 \fulref{p.23}]
\begin{enumerate}[nosep]
\item $\mathrm{Rect}(S)$ 是唯一與 $w(S)$ 等價的表；$\mathrm{Rect}(S)=\mathrm{Rect}(S')$ 若且唯若 $w(S)\equiv w(S')$。
\item 乘積 $T\cdot U$ 的三種做法——把 $w(U)$ 插入 $T$、$\mathrm{Rect}(T*U)$、唯一與 $w(T)\,w(U)$ 等價的表——結果相同。
\end{enumerate}
\end{thmbox}

第 1 點由命題 2 得到，所以 $\mathrm{Rect}(S)$ 與挑內角的順序無關（Claim 2）。例如 p.12 的斜表 $S$，
$w(S)=5\,5\,2\,3\,6\,2\,4\,4\,6\,2$，把它插入空表與把 $S$ 滑成表，都得到 $\yts{22246,346,55}$。
第 2 點：三種做法的字都與 $w(T)\,w(U)$ 等價，所以是同一個表（Claim 3）；而字的連接顯然有結合律，所以乘積也有（Claim 1）。
圖~\ref{fig:f-recap} 的例子 $T\cdot U=\yts{11223,2235,345,466,5}$ 三種做法都得到。

\textbf{Plactic monoid} \fulref{p.23--24}。把字依 Knuth 等價分類，等價類的乘法就是把字接起來（等價的字接起來仍等價），
得到一個么半群 $\mathcal{M}$，Lascoux 與 Schützenberger 稱為 plactic monoid。主定理說：$\mathcal{M}$ 和「表＋乘積」是同一個東西。
再取以表為基底的自由 $\mathbb{Z}$-模、以表的乘積為乘法，得到 \textbf{tableau ring} $R_{[m]}$（數字取自 $[m]=\{1,\ldots,m\}$；
有結合律、不交換）。$T\mapsto x^T$（數字 $i$ 出現幾次，$x_i$ 就幾次方）是 $R_{[m]}\to\mathbb{Z}[x_1,\ldots,x_m]$ 的環同態。

% ---------------------------------------------------------------------
\subsection{Pieri 公式}\label{sec:f-pieri}

在 tableau ring 裡，令 $S_\lambda=S_\lambda[m]$ 為所有形狀 $\lambda$、數字在 $[m]$ 的表的和；它的像就是 Schur 多項式 $s_\lambda(x_1,\ldots,x_m)$。

\begin{thmbox}[Pieri 公式 \fulref{p.24}]
\[
\begin{aligned}
S_\lambda\cdot S_{(p)}&=\sum_{\mu} S_\mu &&(\mu/\lambda\ \text{是 $p$ 格的\textbf{水平帶}：任兩格不同行}),\\
S_\lambda\cdot S_{(1^p)}&=\sum_{\mu} S_\mu &&(\mu/\lambda\ \text{是 $p$ 格的\textbf{垂直帶}：任兩格不同列}).
\end{aligned}
\]
\end{thmbox}

理由是第 1 章 \S1.1 的命題（動畫 F05）：依序插入一列（弱遞增的數），新格子一格比一格靠右、而且不會往下，所以新長出的格子是水平帶；
反過來，若 $U$ 的形狀 $\mu$ 比 $\lambda$ 多一條 $p$ 格的水平帶，從最右邊的新格子開始反向擠出 $p$ 次，
就把 $U$ \textbf{唯一地}拆成 $T\cdot V$（$T$ 的形狀 $\lambda$、$V$ 是一列）。插入一行（嚴格遞增的數）的情形同理，得到垂直帶。

\begin{figure}[H]
\centering
\FulPieri
\caption{$\lambda=(2,1)$、$p=2$。加\textcolor{cLeft}{水平帶}得到 $(4,1)$、$(3,2)$、$(3,1,1)$、$(2,2,1)$；
加\textcolor{cRight}{垂直帶}得到 $(3,2)$、$(3,1,1)$、$(2,2,1)$、$(2,1,1,1)$。}
\label{fig:f-pieri}
\end{figure}

數一數 $m=3$ 的情形：形狀 $(2,1)$ 的表有 8 個、長 2 的列有 6 個，$8\cdot 6=48=24+15+6+3$，右邊正是四種 $\mu$ 的表數。
把每個表換成 $x^T$，一列的表的和是完全對稱多項式 $h_p$、一行的表的和是基本對稱多項式 $e_p$，得到
\[
s_\lambda\,h_p=\sum s_\mu\ (\text{水平帶}),\qquad s_\lambda\,e_p=\sum s_\mu\ (\text{垂直帶}).
\]
用三個變數比較 $x_1^2x_2^2x_3$ 的係數：$s_{(2,1)}h_2$ 那邊是 $6=2+2+1+1$，$s_{(2,1)}e_2$ 那邊是 $4=2+1+1+0$
（最後一項是 0：$s_{(2,1,1,1)}$ 要 4 個不同的數，三個變數不夠）。

% ---------------------------------------------------------------------
\subsection{Kostka 數與兩種序}\label{sec:f-kostka}

表的\textbf{內容}是 $\mu=(\mu_1,\mu_2,\ldots)$：數字 $i$ 出現 $\mu_i$ 次。\textbf{Kostka 數} $K_{\lambda\mu}$ = 形狀 $\lambda$、內容 $\mu$ 的表的個數 \fulref{p.25}。
等價的說法：把數字 $\le k$ 的格子取出來，得到形狀 $\lambda^{(k)}$，於是 $\varnothing\subset\lambda^{(1)}\subset\cdots\subset\lambda$，每一步多一條 $\mu_k$ 格的水平帶；
$K_{\lambda\mu}$ 就是這種鏈的個數。

\begin{figure}[H]
\centering
\FulKostkaTabs
\caption{$\lambda=(3,2)$、$\mu=(2,2,1)$ 的兩個表（\textcolor{cLeft}{1}、\textcolor{cB}{2}、\textcolor{cRight}{3} 分色）與它們的水平帶鏈，所以 $K_{(3,2),(2,2,1)}=2$。}
\label{fig:f-kostka-tabs}
\end{figure}

反覆使用 Pieri 公式 $s_\lambda h_p=\sum s_\mu$，從 1 開始依序乘上 $h_{\mu_1},h_{\mu_2},\ldots$，每乘一次就加一條水平帶：
\[
h_{\mu_1}h_{\mu_2}\cdots h_{\mu_\ell}=\sum_\lambda K_{\lambda\mu}\,s_\lambda ,\qquad\text{例如}\quad h_2h_1h_1=s_{(4)}+2s_{(3,1)}+s_{(2,2)}+s_{(2,1,1)} .
\]
同理 $e_{\mu_1}\cdots e_{\mu_\ell}=\sum_\lambda K_{\tilde\lambda\mu}\,s_\lambda$（$\tilde\lambda$ 是共軛分割；垂直帶轉置後就是水平帶）。

\textbf{兩種序} \fulref{p.26}。\textbf{字典序}：從第一項比起，第一個不同的地方較大者為大，這是全序。
\textbf{支配序} $\mu\trianglelefteq\lambda$：每個 $i$ 都有 $\mu_1+\cdots+\mu_i\le\lambda_1+\cdots+\lambda_i$。
包含 $\Rightarrow$ 支配 $\Rightarrow$ 字典序，反過來都不對：$(2,2)\trianglelefteq(3,1)$ 但 $(2,2)\not\subset(3,1)$；
$(4,1,1)$ 的字典序比 $(3,3)$ 大，卻不支配它（$4+1<3+3$）；$(2,2,2)$ 與 $(3,1,1,1)$ 在支配序下不可比。

由定義，$K_{\lambda\lambda}=1$，而 $\mu$ 的字典序比 $\lambda$ 大時 $K_{\lambda\mu}=0$（第一列放不下那麼多 1，依此類推）。
所以把分割依字典序由小到大排，Kostka 矩陣是對角線全為 1 的下三角矩陣（表~\ref{tab:f-kostka}）。
更精確地（練習 2）：$K_{\lambda\mu}\ne 0$ 若且唯若 $\mu\trianglelefteq\lambda$；字典序較大還不夠，例如 $K_{(4,1,1),(3,3)}=0$。

\begin{table}[H]
\centering
\FulKostkaMatrix
\caption{$n=4$ 的 Kostka 矩陣 $K_{\lambda\mu}$（列 $\lambda$、行 $\mu$，字典序由小到大）：下三角、對角線為 1。}
\label{tab:f-kostka}
\end{table}

下三角矩陣可以反解，所以 $s_\lambda$ 能寫成 $h_\mu$ 的整係數組合，而 $h_\mu$ 是對稱多項式——\textbf{Schur 多項式是對稱的}。
這也推出 $K_{\lambda\mu}$ 與 $\mu$ 各項的順序無關（例如 $K_{(3,2),(2,2,1)}=K_{(3,2),(1,2,2)}=K_{(3,2),(2,1,2)}=2$），第 4 章會直接證明（\S\ref{sec:f-apps}）。

% ---------------------------------------------------------------------
\subsection{行字}\label{sec:f-colword}

另一種把表攤平的方法 \fulref{p.27}：由左到右逐行，每一行\textbf{由下往上}讀，得到\textbf{行字} $w_{\mathrm{col}}(T)$。
圖~\ref{fig:f-recap} 的 $T$：$w_{\mathrm{col}}(T)=5421\ 6432\ 652\ 53$，而 $w_{\mathrm{row}}(T)=56\ 446\ 2355\ 1223$。兩個字長得不一樣，但

\begin{thmbox}[(10) \fulref{p.27}]
$w_{\mathrm{col}}(T)\equiv w_{\mathrm{row}}(T)$。（把行字插入空表，長出來的正是 $T$。）
\end{thmbox}

這由一個更一般的事實得到。\textbf{Lemma 1}：斜表被一條水平或垂直的線切成 $T'$、$T''$，則 $w(T)\equiv w(T')\cdot w(T'')$。
水平切是字的定義本身。垂直切：把左塊 $T'$ 往下移到右塊 $T''$ 的左下，得到斜表 $T'*T''$，它的字正是 $w(T')\,w(T'')$；
再從 $T'$ 正上方的內角由右到左一格一格滑回去，就回到 $T$，由命題 2 得證（圖~\ref{fig:f-cut}）。
一次切下最左邊的一行，對剩下的再切，得到 $w(T)\equiv w(C_1)\,w(C_2)\cdots$，而單獨一行的字就是由下往上讀，所以就是 (10)。

\begin{figure}[H]
\centering
\FulCut
\caption{在第 2 行之後垂直切開：$T'$（左兩行）與 $T''$；$T'*T''$ 的字是 $56442312\cdot 65523$，滑回去仍是 $T$。}
\label{fig:f-cut}
\end{figure}

\begin{notebox}
還有更強的 Lemma 2 \fulref{p.28}：$w_{\mathrm{col}}(T)$ 與 $w_{\mathrm{row}}(T)$ 只用 $K'$ 一種變換（及其反方向）就能互變。書到附錄 A 才用到。
\end{notebox}
